
\documentclass[a4paper,12pt]{article}

\input{glyphtounicode}
\pdfgentounicode=1

\usepackage[T1]{fontenc}
\usepackage[utf8]{inputenc}
\usepackage[english]{babel}
\usepackage[version=4]{mhchem}
\usepackage[dvipsnames]{xcolor}
\usepackage[backend=biber,style=ieee]{biblatex}
\usepackage{newtxtext,newtxmath}
\usepackage{microtype}
\usepackage{setspace}
\usepackage{titlesec}
\usepackage{tocloft}
\usepackage{indentfirst}
\usepackage{enumitem}
\usepackage{fancyhdr}
\usepackage{hyperref}
\usepackage{tabularx}
\usepackage{array}
\usepackage{makecell}
\usepackage{booktabs}
\usepackage{graphicx}
\usepackage{amsmath}
\usepackage{footmisc}
\usepackage{csquotes}
\usepackage{fontawesome5}
\usepackage[
  a4paper,
  top        = 2.00cm,
  bottom     = 2.00cm,
  left       = 2.00cm,
  right      = 2.00cm,
  headheight = 15pt,
  headsep    = 5pt,
  footskip   = 15pt
]{geometry}

\newcommand{\degreeawardvalue}{}
\newcommand{\universityvalue}{}
\newcommand{\addressvalue}{}
\newcommand{\unilogovalue}{}
\newcommand{\copyyearvalue}{}
\newcommand{\defenddatevalue}{}
\newcommand{\orcidvalue}{}
\newcommand{\rightsstatementvalue}{}
\newcommand{\degreeaward}[1]{\def\degreeawardvalue{#1}}
\newcommand{\university}[1]{\def\universityvalue{#1}}
\newcommand{\address}[1]{\def\addressvalue{#1}}
\newcommand{\unilogo}[1]{\def\unilogovalue{#1}}
\newcommand{\copyyear}[1]{\def\copyyearvalue{#1}}
\newcommand{\defenddate}[1]{\def\defenddatevalue{#1}}
\newcommand{\orcid}[1]{\def\orcidvalue{#1}}
\newcommand{\rightsstatement}[1]{\def\rightsstatementvalue{#1}}
\newcommand{\titleskip}{\vskip 4\bigskipamount}
\newcommand{\authorskip}{\vskip 2\bigskipamount}
\newcommand{\resumesection}[2]{\section{\textcolor{red}{#1}#2}}
\newcommand{\resumeSubHeadingListStart}{\begin{itemize}[leftmargin=0.0in, label={}, itemsep=2pt]}
\newcommand{\resumeSubHeadingListEnd}{\end{itemize}}
\newcommand{\resumeItemListStart}{\begin{itemize}}
\newcommand{\resumeItemListEnd}{\end{itemize}}
\newcommand{\resumeItem}[1]{
    \item\small{#1}
}
\newcommand{\resumeSubheading}[4]{
    \item
    \begin{tabular*}{1.0\textwidth}[t]{l@{\extracolsep{\fill}}r}
        \textbf{#1} & \textbf{\small #2} \\
        \small#3 & \small #4
    \end{tabular*}
}

\renewcommand{\cftsecleader}{\cftdotfill{\cftdotsep}}
\renewcommand{\cftsecpresnum}{Chapter~}
\renewcommand{\cftsecaftersnum}{:}
\renewcommand{\headrulewidth}{0pt}
\renewcommand{\footrulewidth}{0pt}
\renewcommand\labelitemi{$\vcenter{\hbox{\tiny$\bullet$}}$}
\renewcommand\labelitemii{$\vcenter{\hbox{\tiny$\bullet$}}$}

\pagestyle{fancy}
\fancyhf{}
\fancyfoot[R]{\thepage}
\urlstyle{same}
\addbibresource{references.bib}
\hypersetup{
  colorlinks=true,
  linkcolor=RoyalBlue,
  urlcolor=RoyalBlue,
  citecolor=RoyalBlue,
  anchorcolor=RoyalBlue
}
\settowidth{\cftsecnumwidth}{Chapter 9: }
\titleformat{\section}
  {\normalfont\Large\bfseries}
  {Chapter \thesection:}
  {0.5em}
  {}
% \titleformat{\section}{
%     \scshape\raggedright\large\bfseries
%     }{}{0em}{}[\color{black}\titlerule]
%     \titlespacing{\section}{0pt}{5pt}{5pt}

\setstretch{1.5}
\setlist[itemize]{topsep=4pt, itemsep=3pt, parsep=0pt, partopsep=0pt, leftmargin=*}
\setlist[enumerate]{topsep=4pt, itemsep=3pt, parsep=0pt, partopsep=0pt, leftmargin=*}
\setlength{\footnotesep}{15pt} 
\setlength{\skip\footins}{15pt} 
\setlength{\parskip}{5pt}

\begin{document}

{\centering
  {\scshape \textls[50]{\fontsize{25}{30}\selectfont Reza {\bfseries Aliasgari Renani}}} \\ 
  {\fontsize{10}{10}\selectfont 16/01/2026} $|$
  {\fontsize{10}{10}\selectfont Motivation Letter} $|$
  {\fontsize{10}{10}\selectfont PhD (Physics, Electrical Engineering)} $|$
  {\fontsize{10}{10}\selectfont Technology Development (IV)} \\ 
  {\fontsize{10}{10}\selectfont Karlsruhe Institute of Technology, Karlsruhe
  School of Elementary Particle and Astroparticle Physics (KIT, KSETA)}
\par}
\vspace{-15pt}
\noindent\rule{\textwidth}{0.5pt}

\vspace{-5pt}

My research objective is to develop microelectronic systems for experimental particle physics, including data acquisition and trigger electronics, real-time processing, and architectures based on field-programmable gate arrays (FPGAs) and system-on-chip (SoC) designs. The KSETA doctoral program at KIT provides the experimental infrastructure and technology development environment needed for this work at the Institute for Data Processing and Electronics. I selected technology development as the primary research field and experimental particle physics as the alternative field because of its intersection with my prior work in FPGA development and experimental physics.

I joined Dr.\ Koveshnikov's laboratory in the Russian Academy of Sciences three years ago to study irradiated semiconductor devices. We investigated the effects of electron and gamma irradiation and hydrogen plasma treatments on \ce{SiO2}-based MOS devices, identifying traps by location (insulator, semiconductor, interface) and carrier type (electrons, holes) using electrical characterization techniques. This work resulted in a first-author publication\footnote{R.~Aliasgari~Renani, O.A.~Soltanovich, M.A.~Knyazev, S.V.~Koveshnikov, ``Investigation of low energy electron irradiated \ce{SiO2} based MOS devices by C--V and TSC techniques,'' \textit{Russian Microelectronics}, 2023, DOI: \href{https://doi.org/10.1134/S1063739723600516}{10.1134/S1063739723600516}} and strengthened my understanding of device stability under bias stress and radiation. I also developed optimized MATLAB automation programs for experimental techniques such as TSC, C--V, I--V, I--T, and DLTS measurements, which enabled systematic data acquisition and reduced experimental run-times.

After my undergraduate degree, I was interested in extending my work from basic semiconductor devices to more complex microelectronic systems and their interaction with radiation. I joined the SoC Development Laboratory to learn about programmable hardware and to work on FPGA-based systems. We implemented image signal processing (ISP) algorithms in Verilog using both manual coding and generated HDL code from Simulink models. I developed a fixed-point arithmetic library, built verification testbenches, and used Python for quality control. We evaluated resource utilization after synthesizing the designs on a Xilinx Zybo Z7 board with live camera input and HDMI output. The codebase was maintained under Git-based source control. 

Concurrently, at the Laboratory of Plasma Systems under the supervision of Dr.\ Vasiliev, we conducted research on radiation and plasma effects on FPGA devices. We exposed an FPGA board to 25--60~keV electron beams in low-pressure oxygen, generating plasma and X-rays from a tungsten target while applying thermal cycling. The FPGA output signal was continuously monitored, and preliminary functionality tests were performed during and after exposure.

I would like to work under the supervision of Dr.\ Becker on embedded electronic systems and adaptive reconfigurable architectures, and to learn hardware and software codesign and low-power design methods through implementation and verification. Dr.\ Kopmann investigates data acquisition, trigger systems, real-time monitoring, GPU computing, and data management for high-rate experiments, and with his guidance I can contribute to architecture design, testbenches, and hardware-in-the-loop validation for these systems. I am also interested in the research by Dr.\ Sander on heterogeneous real-time computing architectures, and I would like to contribute to integrating FPGAs with CPU and GPU control software into stable processing chains.

The multidisciplinary nature of research at KSETA is exactly what I am looking for in my career, and I want to use what I have learned in radiation effects on electronics, device characterization, and system verification to advance its research, from design and verification to experimental testing and long-term reliability assessment. My goal is to contribute to detector electronics and processing systems that remain stable in harsh environments while meeting real-time requirements.

\end{document}
